\documentclass[11pt,twoside]{article}
\usepackage{fontspec}
\usepackage{footmisc}
\usepackage{editio}
\usepackage{polyglossia}
\usepackage[paperwidth=17.50cm,paperheight=25.00cm,top=2.20cm,inner=2.65cm,outer=1.65cm,bottom=2.20cm]{geometry}
\setmainlanguage{spanish}
\setmainfont{Alt Standard}[RawFeature=lnum]
\usepackage[noeledsec,nocritical,nofamiliar,noledgroup]{reledmac}
\linenummargin{right}
\setlength{\linenumsep}{-65pt}
\usepackage[breaklinks=true,hidelinks]{hyperref}
%
\title{\fontsize{40}{20}\selectfont\textbf{HOMÉRICA}}
\date{\vspace{40pt} \textsc{Tradvcción fragmentaria de las obras homéricas por}\\\vspace*{14pt}{\huge\textsc{\textbf{Leopoldo Lvgones}}}\\\fontsize{40}{11}\selectfont{\vfill\textbf{\LaTeX}}}
%\author{}
%
%Para que el título del índice esté en smcp.
\usepackage{tocloft}	
\renewcommand{\cftsecleader}{\cftdotfill{\cftdotsep}}
\renewcommand{\cfttoctitlefont}{\bfseries\scshape}
%Salto de línea.
\newsavebox{\dbox}
\newcommand{\dlgo}[2]{\sbox{\dbox}{#1}\parbox[t]{\linewidth}{#1\strut\\[0pt]\hspace*{\wd\dbox}\hspace{\fontdimen2\font}#2\strut}}
%
\newcommand{\ehuno}{\emph{\textbf{E.-H.\textsuperscript{(1)}}}}
\newcommand{\ehdos}{\emph{\textbf{E.-H.\textsuperscript{(2)}}}}
\newcommand{\ehter}{\emph{\textbf{E.-H.\textsuperscript{(3)}}}}
\newcommand{\ehctr}{\emph{\textbf{E.-H.\textsuperscript{(4)}}}}
\newcommand{\neh}{\emph{\textbf{N. E.-H.}}}
%
\newcommand{\separator}{\noindent-- -- -- -- -- --\par}
%
\usepackage{metalogo}
\usepackage{lettrine}
\AtBeginDocument{\setcounter{ballast}{50}}
%
%
\begin{document}
\pagenumbering{gobble}
	\maketitle
\newpage
%
\null
\vfill
\begin{flushright}
	Edición y composición en \XeLaTeX\ por Ioannes Bracmannus (2025).
\end{flushright}
\newpage
%
\tableofcontents
\newpage
%
\pagenumbering{Roman}
\setcounter{page}{1}
\section*{\centering\textsc{Prefacio}}
\addcontentsline{toc}{section}{\textsc{Prefacio}}\par
%
\noindent\textsc{Las traducciones}, en el área de la literatura, además de tener su valor filológico o científico por su precisión, tienen valor estético. El gran \textsc{Leopoldo Lugones}, probablemente el mayor escritor que hubo en la historia de la literatura argentina, tradujo parte de las obras homéricas, desde episodios enteros hasta apenas unas líneas; y tal es su valor estético que merecen ser editadas por lo mismo, como cualquier traducción. El traductor, además de eximio helenista, fue un un auténtico compositor, un genio del \emph{verbo poético}; por lo mismo, además de imponerse la tarea de traducir fielmente a Homero desde su lengua, trasladando las acepciones al castellano con la mayor fidelidad posible, se impuso también la tarea de trasladar el ritmo, la forma. Para la traducción Lugones adopta el metro alejandrino por considerarlo, en su extensión y estructura, el más cercano al hexámetro dactílico griego, y es así que traduce la primera línea como «Canta, diosa, el encono de Aquiles Peleyades», a partir del mítico «Μῆνιν ἄειδε θεὰ Πηληϊάδεω Ἀχιλῆος».\par
%
El sistema de siglado para agilizar las referencias es el siguiente: \ehuno = \emph{Estudios Helénicos I: La funesta Helena} (1923); \ehdos = \emph{Estudios Helénicos II: Un paladín de la Ilíada} (1923); \ehter = \emph{Estudios Helénicos III: La dama de la Odisea} (1924); \ehctr = \emph{Estudios Helénicos IV: Héctor el domador} (1924); \neh = \emph{Nuevos Estudios Helénicos} (1928). Todos fueron editados por la Biblioteca Argentina de Buenas Ediciones Literarias (Babel).\par
\newpage
%
\pagenumbering{gobble}
\section*{\raggedleft\fontsize{30}{14}\textsc{Ilíada}}
\addcontentsline{toc}{section}{\emph{\textsc{Ilíada}}}
\repage
%
\pagenumbering{arabic}
\setcounter{page}{3}
%Canto 1.
\head{\textsc{\emph{Ilíada}, Canto I}}
\section*{\centering\textbf{\textsc{Canto} I}\footnote{\neh, pp. 58-80.}}
\addcontentsline{toc}{subsection}{\textsc{Canto I}}
%
\beginnumbering
\pstart
\lettrine[lines=4, nindent=2pt]{C}{\textsc{anta, diosa, el encono}} de Aquiles Peleyades,\\
que, aciago, a los aqueos causó males sin cuento,\\
y tantas nobles almas de héroes echó al Hades,\\
dando a perros y buitres sus cuerpos de alimento\\
(pues así los designios Zeus se realizaron)\\
desde que por primera vez riñendo, quebraron\\
el Atrida, rey de hombres, y el deífico Aquiles.\par
¿Cuál fué de entre los dioses el que los puso hostiles?\\
El hijo de Latona y Zeus, quien resentido\\
con el rey, al ejército envió maligna peste\\
que a las tropas diezmaba, porque había ofendido\\
Atrida al sacerdote Crises, cuando llegó éste\\
a las veleras naves aqueas, aportando\\
por rescate de su hija multa inmensa, y llevando\\
en su mano las ínfulas del flechador Apolo\\
a un cetro de oro atadas; con que imploró, no sólo\\
a ambos Atridas, jefes de pueblos, sino a cada\\
\dlgo{uno de los aqueos, diciéndoles:}{--Oh Atridas}\\
y soldados aqueos de las grebas pulidas,\\
quieran los dioses desde su olímpica morada\\
que la ciudad de Príamo toméis, y con buen sino\\
volváis a vuestra tierra; mas, dadme a mi hija amada,\\
aceptando el rescate, y acatando al divino\\
\dlgo{flechero Apolo.}{Todos aclaman que se deba}\\
respeto al sacerdote, recibiendo la multa.\\
Agamenón Atrida solamente no aprueba;\\
antes, con malos modos, lo echa de allá y lo insulta:\par
--Que no vuelva yo, anciano, ni ahora ni luego a hallarte,\\
ya te quedes o tornes, junto a la flota anclada;\\
pues ni el cetro ni la ínfula del dios podrán guardarte.\\
No soltaré a tu hija, antes que envejezca expatriada,\\
tejiendo en casa, en Argos, y durmiendo conmigo.\\
Y vete, y no me irrites, y evita así el castigo.\par
Dice, y medroso el viejo, lo acata y abandona\\
en silencio, y costeando va el estruendoso mar;\\
y al rey Apolo, vástago de la crespa Latona,\\
mientras se aparta, empieza con fervor a rogar:\par
--Oyeme, dios del arco de plata, protector\\
de Crisa y de la santa Kila, dominador\\
de Ténedos, ¡oh Esmintio!: Si tu iglesia graciosa\\
labrar supe, y quemarte pingües muslos de cabras\\
y toros, que los dánaos, conforme a estas palabras,\\
mis lágrimas expíen con tu flecha luctuosa.\par
Rogó así, y Febo Apolo que lo oyó, con airado\\
corazón, baja al punto del Olimpo encumbrado,\\
al hombro arco y aljaba de bien cerrado broche.\\
Y en la espalda resuénanle, al moverse enojado,\\
las flechas, cuando avanza semejante a la noche.\\
Siéntese después, lejos de las naves, y tira\\
una flecha, al terrible son del arco de plata.\\
Mulas y ágiles perros hiere primero su ira;\\
mas, luego, con fatales tiros los hombres mata,\\
y arden sin tregua innúmeros difuntos en la pira.\par
Nueve días las tropas flechó el dios por doquiera;\\
mas, el décimo, en junta convocó al pueblo Aquiles;\\
pues lo inspiró la diosa de blancos brazos, Hera,\\
ansiosa por los dánaos que veía morir.\\
Y luego que, llamados, la junta se reuniera,\\
alzóse el raudo Aquiles, y así empezó a decir:\par
--Atrida, creo que ahora, nuevamente vagando,\\
retroceder debemos, si, la muerte evitando,\\
guerra y peste no rinden juntas a los aqueos.\\
Mas, ea, consultémonos ya de algún adivino\\
sacerdote o intérprete de sueños (pues divino\\
el sueño es también) con que nos diga él sin rodeos\\
qué irrita a Febo Apolo: si nos reclama, acaso,\\
promesas o hecatombes, o si, aceptando el craso\\
humo de albos corderos y de escogidas cabras,\\
\dlgo{nos quitará el plageño.}{Dichas estas palabras,}\\
volvió a sentarse. Entonces se alzó el más excelente\\
augur, Calcas Testórides, quien sabiendo lo que era,\\
es y será, las naves aqueas condujera\\
hasta Ilión, con ayuda del don clarividente\\
que Febo Apolo hiciérale. Y así, benevolente\\
\dlgo{los arengó:}{--Oh Aquiles, caro a Zeus, pues mandarme}\\
quieres que explique la ira de Apolo el real flechero,\\
lo diré; mas prométeme y júrame primero\\
que de palabra y obra socorro has de prestarme;\\
no irrite yo a quien todos los argivos comanda\\
con gran poder, y acatan los aqueos; pues si anda\\
mal con el rey un súbdito como aquél puede más,\\
aunque trague su cólera en el mismo momento\\
hasta satisfacerla guarda el resentimiento\\
en su pecho. Declara, pues, si me salvarás.\par
\dlgo{Contestóle así el rápido Aquiles:}{--Ten confianza,}\\
y el augurio que sabes, dínoslo sin tardanza;\\
pues juro por Apolo, caro a Zeus, y a quien oras,\\
Calcas, cuando a los dánaos su oráculo asesoras,\\
que ni uno de ellos, mientras esté yo vivo y sano,\\
junto a las hondas naves pondrá en ti osada mano;\\
aunque te refieras a Agamenón, que a todos\\
los aqueos se alaba de recuperar ufano.\par
El adivino ilustre tomó confianza y dijo:\\
--No se queja él por votos ni hecatombes, de fijo;\\
mas por su sacerdote que Agamenón vejó,\\
pues no libertó a su hija ni el rescate aceptó.\\
Por esto el Flechero hízonos el mal y todavía\\
nos lo hará, y de las Parcas de la peste sombría\\
no ha de librarnos mientras no quede reparada\\
la ofensa, devolviéndose, sin multa rescatada,\\
al caro padre aquella joven de lindos ojos;\\
y una sacra hecatombe sea en Crisa inmolada.\\
Así conseguiremos aplacar sus enojos.\par
Dichas estas palabras, volvió a sentarse, cuando\\
poseído de angustia se alzó el héroe Atrida,\\
Agamenón potente, la negra entraña henchida\\
de rencor, y con ojos centelleantes mirando\\
\dlgo{a Calcas torvamente, díjole:}{--Mago avieso,}\\
nunca me anuncias nada feliz; antes gozando\\
en predecir desgracias, sólo haces mal con eso.\\
Y ahora entre los dánaos declamas, augurando\\
que el Flechero tortúralos así porque no quise\\
el rescate halagüeño de la doncella Crise,\\
a quien deseaba en casa tener, pues la prefiero\\
realmente a Clitemnestra mi legítima esposa,\\
ya que en cuerpo ni en talle, ni por diestra o graciosa\\
le es inferior; mas, siempre que mejor sea, quiero\\
devolverla, y que el pueblo se salve y no perezca.\\
Pero aprontadme ya otro premio, pues no es cumplido\\
que así el único argivo que ninguno merezca\\
sea yo, viendo todos que el mío habré perdido.\par
Mas, el divino Aquiles respondió de este modo:\\
--Atrida, el más glorioso y avariento de todos,\\
¿qué premios los magnánimos aqueos han de darte,\\
si riquezas comunes no hay en ninguna parte?\\
Repartimos cuanto en las ciudades conquistamos,\\
y no conviene al pueblo que a juntarlo volvamos.\\
Remite al dios la joven, y cuando a los aqueos\\
Zeus nos entregue Troya, la del fuerte bastión,\\
triple o cuádruple pago te asignará el saqueo.\par
Respondió a estas palabras el rey Agamenón:\\
--Por más bravo que seas, no intentes con desvío\\
traerme, ínclito Aquiles, a engaño o persuasión.\\
¿Por conservar tu premio, quieres que deje el mío,\\
y para esto me impones que entre la doncella?\\
Lo haré si los magnánimos aqueos dan por ella\\
algo que la compense, conforme a mi albedrío.\\
Y si no me lo dieran, yo mismo iré a tomar\\
tu premio, o el de Ulises, o el de Ayax, y el que sea\\
se indignará. Mas, de esto, después hemos de hablar.\\
Ahora, pronto botemos a la divina mar\\
negro bajel que una hábil tripulación provea;\\
y una hecatombe, junto con la lozana Crise,\\
embarquemos, poniéndolas al mando de un guerrero:\\
Ayax, Idomeneo, o el deífico Ulises,\\
o tú, Pelida, que eres de todos el más fiero,\\
para que con las víctimas aplaques al Flechero.\par
Mas, de través mirándolo, contestó el raudo Aquiles:\\
--Ay de mí, avaro impúdico ¿habrá un argivo, acaso,\\
que obedezca tus órdenes, ya siguiendo tus pasos,\\
ya afrontando, valiente, los combates viriles?\\
No fue por los lanceros troyanos provocado\\
como vine a la lucha, que en nada me ofendieron,\\
ni jamás mis caballos y vacas se han llevado,\\
ni en Ptía la prolífica y feraz, destruyeron\\
nuestras cosechas, porque siempre nos dividieron\\
muchas sierras umbrías y el mar alborotado.\\
Mas, contigo vinimos, oh infame, a que gozaras\\
con dar en los de Troya venganza a Menelao,\\
y a ti, cara de perro, que ya en nada reparas.\\
Y amenazas privarme del premio que me dieron\\
por mis grandes fatigas los dánaos, aunque poca\\
sale siempre ante el tuyo la parte que me toca\\
de las ricas ciudades troyanas que rindieron.\\
Si bien yo aguanto el peso más arduo de la lucha,\\
y al hacerse el reparto tu recompensa es mucha\\
y la mía pequeña, siempre a mis naves vuelvo\\
contento, aunque cansado de pelear. Mas, resuelvo\\
tornar ya a Ptía, viendo que irme a casa es mejor\\
en las cóncavas naves, que a costa de mi honor,\\
lucro y caudal seguirte granjeando aquí a la vez.\par
\dlgo{Agamenón, rey de hombres, repúsole:}{--Huye, pues,}\\
si se te antoja. No te ruego que a causa mía\\
te quedes. Me harán otros honrosa compañía,\\
y sobre todo el próvido Zeus. Que tú, el más odioso\\
me eres entre los reyes a quienes Dios ayuda,\\
pues de riña y discordias andas siempre gustoso.\\
Y si tu fuerza es tanta, débeslo a un dios, sin duda.\\
Vete con naves y hombres, y reina en paz completa\\
sobre los mirmidones, que tu ira no me inquieta.\\
Mas ya que Febo Apolo va Criseida a quitarme,\\
y con mi buque y gente debo yo enviarla en prenda\\
te amenazo con irme yo en persona a tu tienda,\\
y a la linda Briseida, tu galardón, llevarme;\\
para que veas cómo soy yo el que puede más,\\
y nadie ose conmigo compararse jamás.\par
Tal dijo, y el Peliónida sintió que lo angustiaba\\
su corazón, dudando dentro el pecho velludo,\\
si sacaba a lo largo del muslo el sable agudo\\
y apartando a los otros, al Atrida mataba,\\
o si su furia y cólera reprimía y calmaba.\\
Mientras así agitábanse su corazón y mente,\\
y que desenvainando ya iba la espada ingente,\\
bajó Atenea del cielo, pues la enviaba la diosa\\
de blancos brazos, Hera, que amaba cordialmente\\
a los dos, y por ambos andaba cuidadosa.\\
De atrás tiró al Pelióńida por los cabellos rojos,\\
visible a él solo, que otro ninguno la ha advertido;\\
volvióse aquél atónito, y al punto ha conocido\\
a Palas Atenea por sus terribles ojos;\\
con que aladas palabras así le ha dirigido:\par
--Hija del dios portaégida, ¿a qué viniste? ¿Acaso\\
a presenciar la ofensa de Agamenón Atrida?\\
Mas, puedo predecirte como seguro caso,\\
que pronto su insolencia le costará la vida.\par
Atena la ojizarca contestóle prudente:\\
--He venido del cielo para hacer que, obediente,\\
tu cólera reprimas; pues me envía la diosa\\
de blancos brazos, Hera, que os ama cordialmente\\
a los dos, y que andaba por ambos cuidadosa.\\
Cesa, pues, la disputa, no desnudes la espada,\\
y oféndelo de boca, si hacerlo así te agrada.\\
Te prometo que un día, triple y brillante premio\\
tendrás por este ultraje. Mas, cede a nuestro apremio\\
\dlgo{Y acátanos ya.}{Aquiles de raudos pies contesta:}\\
--Más vale observar, diosa, vuestra orden sin protesta,\\
aunque el corazón tengo grandemente enojado;\\
pues a quien oye a los dioses, es de ellos escuchado.\par
Dijo, y puesta la recia mano al puño argentino,\\
envainó la ancha espada, dócil a Atena, y ella\\
tornó al Olimpo donde con los otros divinos\\
\dlgo{Zeus portaégida mora.}{Mas, con nueva querella,}\\
al Atrida el peliónida increpó siempre acerbo:\par
--Borracho de ojos cínicos y corazón de ciervo,\\
nunca osaste en la guerra con los del pueblo armarte\\
o ir con los principales aqueos al acecho,\\
por miedo a la muerte. Sin duda es más provecho,\\
en el gran campo aqueo despojar de su parte\\
a quien te contradiga. Rey que voraz impera\\
sobre menguada gente, porque de otra manera\\
lanzabas hoy, Atrida, tu último insulto aquí,\\
sobre gran prenda quiero jurarte algo más: ¡Sí,\\
por este cetro que hojas ni ramas dará nunca,\\
pues desde que en el monte dejó su cepa trunca,\\
el bronce al tornearlo privólo, con certeza,\\
de que le retoñaran el follaje ni corteza;\\
y que en sus manos llevan los dánaos al presente,\\
cuando según las leyes de Dios juzgan fielmente:\\
con tal prenda te juro que un día todo aqueo\\
se acordará de Aquiles, y en tu vano deseo\\
de ayudarlos, cuando Héctor tantos mate valiente,\\
desgarrará honda pena en tu corazón furioso,\\
porque al principal de ellos no honraste respetuoso!\par
Tal se expresó el Pelida, contra el suelo arrojando\\
su cetro tachonado de oro, y sentóse enfrente\\
del Atrida que se iba con más furor cargando.\\
Pero entonces de entre ellos se alzó el blandilocuente\\
orador de los pilios, Néstor, cuya cadente\\
palabra, de su lengua, como la miel fluía.\\
Ya dos generaciones de hombres de artuculada\\
voz, que fueron nacidas y consigo criadas\\
en la divina Pylos, ante él pasado habían.\\
Y sobre la tercera mandaba todavía.\\
Este fué el que, arengándolos, benevolente dijo:\par
--¡Oh dioses, qué gran duelo para la tierra aquea!\\
Y cómo a los troyanos, y a Príamo y sus hijos,\\
satisfacción inmensa causaría, de fijo,\\
saber vuestro altercado, siendo que en la pelea\\
y el consejo estáis sobre todos los dánaos. Pero\\
desde que sois más jóvenes, condescended primero;\\
pues yo frecuenté otrora varones principales,\\
mejores que nosotros, y no me desoyeron.\\
Todavía no he visto ni veré hombres iguales\\
a Piritoo, Driante rey de pueblos, Keneo,\\
Exadio y el deífico Polifemo, y Teseo\\
el Egeída, que era como los inmortales.\\
Bravos entre los bravos que alimentó la tierra,\\
ellos a los centauros monteses dieron guerra\\
por más bravos, matándolos en terrible exterminio.\\
Para incorporarme a ellos, dejé la tierra Apía,\\
pues fuí allá desde Pylos, mi lejano dominio,\\
y a su propio llamado luché por cuenta mía.\\
Ningún hombre de ahora resistirlos podría;\\
y, no obstante, acataban mis consejos realmente.\\
Prestadme, así, obediencia, que es lo más conveniente.\\
Ni tú, aunque el mejor seas, la muchacha le quites,\\
puesto que los aqueos con ella lo han premiado.\\
Ni tú, Pelida, quieras altercar, igualado\\
con el rey, ya que en honra ninguno le compite\\
de entre los portacetros por Zeus glorificados.\\
Ahora, si por tu parte tienes más valentía,\\
y si diosa es la excelsa madre que a luz te ha dado,\\
el puede más, pues reina sobre la mayoría.\\
Y tú, Atrida, apacigua tu cólera, y depón\\
tu encono contra Aquiles, que gran resguardo presta\\
a todos los aqueos en la guerra funesta.\par
Contestó, así diciéndole, el rey Agamenón:\\
--Sí, cuanto has dicho, anciano, por cierto es de razón;\\
pero este hombre pretende sobreponerse a todos,\\
y a todos dirigirlos y mandar, de tal modo\\
que nadie ha de aceptárselo. ¿Por que bravo lo hicieron\\
los dioses, que insultara también le permitieron?\par
Mas el divino Aquiles lo interrumpió exclamando:\\
--Por vil y por cobarde con razón pasaría,\\
si a todo cuanto dices cediera. Impón el mando\\
a otros, no a mí, pues nunca más pienso obedecerte.\\
Pero oye esto y consérvalo en tu alma bien grabado:\\
no disputaré a nadie ni a ti, con mano fuerte,\\
la joven, ya que en suma quitáis lo que habéis dado;\\
mas, de lo que me resta junto al negro y ligero\\
bajel, nada podrías tomar, si yo no quiero.\\
De lo contrario, inténtalo; y así éstos sin tardanza\\
verán tu negra sangre chorreando por mi lanza.\par
Habiendo altercado ambos así, en pie se pusieron,\\
y ante la flota aquea la junta disolvieron.\\
El Pelida, a las tiendas y parejos bajeles\\
se fué con Meneciades y sus amigos fieles.\\
Y en tanto, al mar, Atrida, botó un barco ligero;\\
eligió por su propia cuenta veinte remeros,\\
cargó allá una hecatombe para el dios, y llevando\\
a la linda Criseida, la embarcó con esmero.\\
Y el ingenioso Ulises fué quien tomó el comando.\par
Tras que así por la líquida ruta a bogar echaran,\\
mandó Atrida a los pueblos que se purificaran,\\
cual lo hicieron, lanzando toda inmundicia al mar.\\
Luego a Apolo inmolaron en la playa, perfectas\\
hecatombes de toros y de cabras selectas;\\
y el olor iba al cielo con el humo al ondear.\par
Agamenón, en tanto que esto hacía la gente,\\
sin olvidar que a Aquiles en el agrio debate\\
amenazó primero, dijo a sus diligentes\\
heraldos y criados, Taltibio y Euribate:\par
--A la tienda de Aquiles Peleyades id presto,\\
y a la linda Briseida traedme de la mano;\\
o he de ir allá en persona, si él no la da de plano,\\
con más gente a sacársela, y peor le saldrá esto.\par
Despachólos, vertiendo tan fieras expresiones;\\
y ellos, el mar costeando, con ánimo indispuesto,\\
llegaron a las tiendas y barcos mirmidones.\\
Junto a su tienda y negro buque, sentado hallaron\\
a Aquiles, quien, de verlos, no se alegró por cierto;\\
y haciendo al rey la venia, con grave desconcierto,\\
parados ante él nada dijéronle ni hablaron.\\
Pero él entendió en su alma, y alzando así la voz,\\
\dlgo{dijo:}{--Salud, heraldos, mensajeros de Dios}\\
y de los hombres. Aproximaos sin cuidado,\\
pues la culpa no es vuestra, sino de Agamenón,\\
que en busca de la joven Briseida os ha mandado.\\
Vamos, noble Patroclo, saca a disposición\\
de ambos la joven que han de llevarse; y que testigos\\
me sean los dos ante los dioses venturosos,\\
y los hombres mortales, y ese rey rencoroso,\\
si contra el vil desastre cuentan después conmigo.\\
Pues aquél, poseído de furor, ya no sabe\\
ni prever ni acordarse, para que así, al abrigo,\\
combatan los aqueos con él junto a las naves.\par
Dijo, y sacó Patroclo, dócil al caro amigo,\\
de la tienda a la linda Briseida, y ya entregada,\\
tornaron los heraldos, que siguió contrariada\\
la mujer, la aquea flota. Rompiendo en llanto,\\
Aquiles, de los suyos, lejos se apartó en tanto,\\
y sentado a la orilla del espumoso mar,\\
frente al Océano obscuro, con las manos tendidas,\\
a su querida madre comenzó así a implorar:\par
--Madre, ya que me diste tú a luz de corta vida,\\
pudo Zeus, el olímpico altitonante, honrarme.\\
Mas nada hizo, y el déspota Agamenón Atrida,\\
me ha ultrajado, llegando del premio a despojarme.\par
Así exclamó llorando. Lo oyó su augusta madre,\\
que sentada se hallaba junto al anciano padre\\
en el fondo marino, y al punto de la espuma\\
se alzó como una niebla; y ante el lloroso hijo\\
sentóse, acariciándolo con la mano, y le dijo:\par
--Hijo, por qué llorabas? Qué pena tu alma abruma?\\
Habla, y sepámoslo ambos. Nada oculte tu mente.\par
Aquiles contestóle, suspirando hondamente:\\
--Lo sabes. ¿A qué habría de hablar si sabes todo?\\
A Tebas la sagrada ciudad de Eecion fuimos,\\
y habiéndola saqueado, todo acá nos trajimos.\\
Los aqueos partiéronlo de equitativo modo,\\
y eligieron la linda Criseida para Atrida.\\
Mas Crises, sacerdote de Apolo, fué en seguida\\
a las veleras naves aqueas, aportando\\
por rescate de su hija multa inmensa, y llevando\\
en su mano las ínfulas del flechador Apolo\\
a un cetro de oro atadas; con que imploró, no sólo\\
a ambos Atridas, jefes de pueblos, sino a cada\\
\dlgo{uno de los aqueos.}{Fué de éstos aclamado,}\\
respeto al sacerdote, y el rescate aceptado.\\
Mas no plugo al Atrida la idea así acordada,\\
por lo cual, insultándolo, echó a aquél malamente.\\
Fuése el viejo irritado, y Apolo oyó su ruego,\\
que mucho lo quería; con que tiró inclemente\\
saeta a los argivos, cuyos soldados luego\\
morían a montones; pues las divinas flechas,\\
por todo el grande ejército Aqueo hacían brechas.\\
Díjonos los oráculos del Flechero un sapiente\\
adivino, y consejo di yo inmediatamente\\
de apaciguar al numen. Entonces el Atrida,\\
encendiéndose en cólera, alzóse bruscamente,\\
y me echó una amenaza que está, en verdad, cumplida.\\
Pues ya los perspicaces dánaos a Crise enviaron\\
junto con reales dones, en nave bogadora;\\
mientras que mi tienda unos heraldos se llevaron\\
a la hija de Briseo que aquéllos me otorgaron.\\
Si a tu hijo auxiliar quieres, ve al Olimpo e implora\\
a Zeus, que acaso un día su corazón tocaste\\
con hechos o palabras; porque recuerdo ahora,\\
que en el paterno alcázar mil veces te gloriaste\\
de que sola entre todos los inmortales fueras\\
quien al Kronión sombrío vergonzoso contraste\\
supo evitar, cuando otros Olímpicos, con Hera\\
y Poseidón, y Palas Atenea, qusieran\\
encadenarlo. Pero tú, diosa, allí acudiste,\\
desataste sus lazos, y al punto condujiste\\
al vasto Olimpo el fiero centímano que llaman\\
Briareo los dioses y Egeón los humanos.\\
Más fuerte que su padre, fué, de su gloria ufano,\\
a sentarse a la vera del Kronión; y medrosos\\
de él, ya al otro no ataron los dioses venturosos.\\
Recuérdaselo ahora, y llégate a él, y abrazando\\
sus rodillas, consigue que a los troyanos dando\\
su ayuda, ante la flota rechacen y destrocen\\
estos a los aqueos, que así de su rey gocen.\\
Y su error vea el déspota Agamenón Atrida,\\
porque al principal de ellos no honró en forma cumplida.\par
\dlgo{Respondióle llorando Tetis:}{--¿Por qué, hijo mío,}\\
si te di a luz en hora fatal, te habré criado?\\
Sin lágrimas ni pena debiste en los navíos\\
quedar, ya que a tan próximo fin estás destinado.\\
Y ahora eres corto en días y el más desventurado;\\
que al ser en el palacio te di con hado impío.\\
A hablar por ti al fulmíneo Zeus iré yo al nevado\\
Olimpo, y quizá pueda persuadirlo. Y ahora,\\
con los dánaos airado, gana tu bogadora\\
flota, y no hagas más guerras. Hacia el Océano ayer\\
partió Zeus a un banquete que dan los caballeros\\
etíopes y todos los dioses con él fueron.\\
Mas al Olimpo en doce días debe volver;\\
y a su alcázar de bronce, yendo en tu nombre a hablarle,\\
asiré sus rodillas y espero propiciarle.\par
Dijo y partió, dejándolo con el alma enconada,\\
por la elegante esposa que le fué arrebatada.\\
Y mientras tanto Ulises a Crisa conducía\\
la hecatombe sagrada. Tan luego que arribaron\\
al puerto de hondas aguas, amainando guardaron\\
las velas en el negro buque; y en la crujía,\\
al mástil, prontamente, con cables abatieron.\\
Bogaron a remo hasta la rada donde echaron\\
anclas; y atando amarras, a la costa bajaron.\\
sacaron la hecatomble para Apolo el flechero,\\
y al fin salió Criseida del bajel marinero.\\
El ingenioso Ulises condújola al altar,\\
y entregándola al caro padre empezó así a hablar:\par
--Oh Crises, el rey de hombres Agamenón me envía\\
a que tu hija te traiga, y a que en ofrenda pía\\
inmole yo a aquí a Febo la hecatombe sagrada\\
en nombre de los dánaos, a ver si el dios se apiada,\\
pues mandó a los argivos lamentable avería.\par
Dijo, y su hija entrególe, que él recibió contento.\\
Y al punto, en torno al ara de sólido cimiento,\\
para el dios ordenaron la hecatombe sagrada.\\
Laváronse las manos, tomaron la salada\\
harina, y allá entonces, en alta voz orando,\\
Crises rogó por ellos, las manos levantando:\par
--Oyeme, dios del areo de plata, protector\\
de Crisa y de la santa Kila, dominador\\
de Ténedos: si un día mis ruegos escuchando\\
me honraste, al pueblo aqueo cruelmente castigando,\\
ahora como entonces, cúmpleme también este voto,\\
y libra a los dánaos de la horrorosa peste.\par
Dijo así, y fué de Apolo su oración escuchada.\\
Luego que oran y esparcen ya la harina salada,\\
mancornando a las víctimas, degüellan y desuellan.\\
Cortan los muslos, cúbrenlos con un doble haz de pella\\
y magras, y el anciano sobre la lecha hacheada\\
quemólos, y con negro vino los fué rociando.\\
En torno suyo había mancebos que empuñando\\
asadores de cinco puntas, se dispusieron.\\
Consumidos los muslos, las entrañas comieron;\\
y ensartando en los pinchos lo restante, lo asaron\\
con cuidado, por trozos; sacáronlo, y cumplida\\
tal faena, el dispuesto banquete comenzaron\\
sin que a nadie faltara su porción. En seguida\\
que saciaron sed y hambre, los mancebos colmaron\\
las cráteras de vino que después repartían,\\
y a todos las primicias de la copa obsequiaron.\\
Los jóvenes aqueos, tras de holgar todo el día,\\
al dios se propiciaron cantándole un hermoso\\
himno, que oyó el Flechero con ánimo gozoso.\par
Cuando el sol se entró y vino la sombra, se acostaron\\
junto al cable del buque; y al asomar temprana\\
la Aurora rosodáctila, hija de la mañana,\\
para el campo aqueo, mar afuera tomaron.\\
Y el flechador Apolo buen viento les envía.\\
Izando luego el mástil, las velas desplegaron,\\
que en blanca combadura con el soplo se hincharon,\\
mientras la honda purpúrea resonando batía\\
la quilla del navío que su rumbo corría.\\
Ya ante el gran campo aqueo, la nave a tierra halaron;\\
y empinándola en grandes vigas sobre la arena,\\
por los buques y tiendas luego se dispersaron.\par
Junto a su flota en tanto, con su alma de ira llena,\\
seguía el raudo Aquiles, noble hijo de Peleo;\\
y de la ilustre junta y el combate alejado,\\
el corazón ya le iba devorando el deseo\\
\dlgo{de griterió y Guerra.}{Mas, cuando hubo llegado}\\
la duodécima aurora, y al Olimpo volvieron\\
los sempiternos dioses por Zeus encabezados,\\
Tetis, que los encargos de su hijo no ha olvidado,\\
emergiendo temprana de las olas del mar,\\
y cruzando el gran cielo, fué en el Olimpo a hallar\\
al vidente Krónida, que se hallaba sentado\\
en el pico más alto del Olimpo escarpado.\\
Fué, pues, y ante él sentándose, las rodillas le asió\\
con la izquierda, y tomándole el mentón con la diestra,\\
su ruego al soberano dios Kronión dirigió:\par
--Padre Zeus, si he dado entre los inmortales muestra\\
de serte útil con actos o palabras, te pido\\
que oigas mi ruego y honres a mi hijo, el de más corta\\
vida entre todos, porque de Agamenón soporta\\
la ofensa que arrancándole su premio le ha inferido.\\
Véngalo, pues, Olímpico Zeus próvido, acordando\\
el triunfo a los troyanos, hasta que los aqueos\\
cumplan con mi hijo y lo honren conforme a sus deseos.\par
Dijo, mas Zeus nubígero, largo tiempo callando,\\
nada respondió. Tetis, que abrazando seguía\\
sus rodillas con fuerza, suplicó todavía:\par
--Prométeme y con una señal veraz asiente,\\
o niega, ya que a nadie temes; y así enterada,\\
sabré si entre las diosas soy la más despreciada.\par
El nubígero Zeus, suspirando hondamente,\\
\dlgo{díjole:}{--Algo funesto se apronta, pues con Hera}\\
me enconarás, cuando ella me provoque y zahiera.\\
Ya entre los dioses ríñeme sin motivo, a pretexto\\
de que yo en la pelea protejo a los troyanos.\\
Pero vete y apártate, no te advierta Hera; y a esto,\\
para que bien se cumpla, déjalo ya en mis manos.\\
Te haré con la cabeza la seña en que confías.\\
Entre los inmortales no hay mayor garantía.\\
Que irrevocable y cierto, de cumplirse no deja\\
lo que con la cabeza ratifiqué a fe mía.\par
Dijo el Kronión. Al signo de las azules cejas,\\
ondeó el fragante pelo del rey en la inmortal\\
cabeza, y el Olimpo retembló colosal.\par
Dicho aquello, apartáronse. Regresó ella a su centro,\\
al hondo mar saltando del Olimpo esplendente,\\
y Zeus a su palacio. Los dioses juntamente\\
ante su padre alzáronse, saliéndole al encuentro.\\
Sentóse él en su trono; mas Hera que con tino\\
había visto cómo con él deliberó\\
Tetis pies de plata, hija del anciano marino,\\
al dios Kronión hirientes palabras dirigió:\par
--¿Cuál de los dioses, pérfido, contigo a conversado?\\
Siempre tramar secretos sin mí, te causa agrado;\\
mas, de lo que meditas, nunca amable me enteras.\par
\dlgo{Contestó el padre de hombres y dioses:}{--En vano, Hera,}\\
pretendes de tal modo saber todas mis cosas,\\
porque arduas te saldrían aunque seas mi esposa.\\
Lo que es de oírse, ni hombre ni dios sabrá primero;\\
mas todo cuanto sin los dioses resolver quiero,\\
\dlgo{No indagues ni preguntes.}{Respondió la augusta Hera}\\
\dlgo{de ojos de buey:}{--¡Terrible Kronida, qué dijiste!}\\
Poco indago o pregunto, puesto que deliberas\\
muy tranquilo tus cosas. Mas tengo el alma triste\\
de recelar que acaso te sedujo con tino\\
Tetis pies dep plata, hija del anciano marino.\\
Temprano vino a asirte las rodillas, y creo\\
que para honrar a Aquiles le acordaste, de fijo,\\
perder a muchos ante los bajeles aqueos.\par
El nubígero Zeus, respondiéndole dijo:\\
--Desgraciada! Sospechas siempre lo que no oculto.\\
Así obtendrás tan sólo quedar más alejada\\
de mi afecto, y con ello sufrir mayor insulto.\\
Si es verdad lo que dices, será porque me agrada.\\
Pero, acata mis órdenes y siéntate callada.\\
Teme que ni los dioses del Olimpi, a ti amables,\\
te valgan si en ti pongo mis manos formidables.\par
Dijo así. La augusta Hera de ojos de buey, temiendo,\\
fué a sentarse callada, su furor reprimiendo.\\
Y gimieron los célicos dioses en la morada\\
de Zeus. Mas ya el artífice Hefesto los arenga\\
mimando a Hera de blancos brazos, su madre amada:\par
--Temo que algo insufrible y aciago sobrevenga,\\
si así por los mortales reñís, alborotando\\
a los dioses, porque hasta se nos irá amargando\\
el placer de las fiestas. Yo aconsejo a mi madre,\\
aun cuando es tan prudente, mime a Zeus, caro padre;\\
con que así él nuestras fiestas no turbe rencilloso.\\
Pues si viene al Olímpico fulminador la idea,\\
puede echarnos abajo, siendo el más poderoso.\\
Hazle en dulces palabras halago cariñoso,\\
a fin de que el Olímpico favorable nos sea.\par
Habló así; y dirigiéndose a su madre querida,\\
puso en sus manos una doble copa y le dijo:\\
--Sufre todo, madre, aunque sea cruel tu quebranto.\\
No te vea golpeada yo, queriéndote tanto;\\
y por más que en su furia nada te valga tu hijo,\\
pues arduo es al Olímpico resistir. Ya una vez\\
que intenté socorrerte, me asió de un pie y después\\
desde el umbral divino me arrojó. Todo el día\\
rodé, cayendo en Lemnos cuando el sol se ponía.\\
Débil soplo quedábame de vida postrimera,\\
y los sintios salváronme del golpe prontamente.\par
Así dijo, y la diosa de blancos brazos, Hera,\\
la copa de su mano recibió sonriente.\\
Luego él, por la derecha sirvió a los otros dioses,\\
sacando de una crátera el dulce néctar. A esto,\\
de los felices númenes inextinguibles alzóse\\
la risa en el palacio, viendo afanarse a Hefesto\par
Hasta que el sol se puso, todo el día gozaron\\
del festín, sin que nadie justa porción faltara,\\
ni la ira magnífica de Apolo, ni la clara\\
voz con que allá las musas alternando cantaron.\par
Mas, luego de apagada la solar refulgencia,\\
ganaron las mansiones que a cada uno había hecho\\
Hefesto el cojo ilustre con ingeniosa ciencia.\\
Zeus, el fulmíneo Olímpico, fué a buscar en su lecho\\
la calma y dulce sueño que hallaba satisfecho.\\
Habiendo allá subido, durmióse; y a su vera,\\
reposaba acostada la crisótrona Hera.
\pend
\endnumbering
\repage
%
%Canto 2.
\head{\textsc{\emph{Ilíada}, Canto II}}
\section*{\centering\textbf{\textsc{Canto} II}\footnote{\neh, pp. 92-93.}}
\addcontentsline{toc}{subsection}{\textsc{Canto II}}
%
\noindent-- -- -- -- -- --\par
%
\beginnumbering
\pstart
\setline{474}Tal como los cabreros distinguen fácilmente\\
las copiosas majadas, si en el pasto se mezclan,\\
así les dan los jefes el orden conveniente\\
para el combate. Entre ellos va el rey Agamenón,\\
parecido por ojos y testa al dios tonante;\\
por el talle a Ares y por el pecho a Poseidón.\\
Cual la mayor cabeza del hato es el pujante\\
toro que entre las vacas sobresale, aquel día\\
dispuso Zeus que Atrida fuera el más arrogante\\
de porte entre los muchos héroes que allí había.\par
\pend
\endnumbering
%
\noindent-- -- -- -- -- --\par
\repage
%
%Canto 3.
\head{\textsc{\emph{Ilíada}, Canto III}}
\section*{\centering\textbf{\textsc{Canto} III}\footnote{\ehuno, pp. 43-46, 49, 52-53, 62 + \ehdos, p. 81 + \neh, p. 24.}}
\addcontentsline{toc}{subsection}{\textsc{Canto III}}
%
\noindent-- -- -- -- -- --\par
%
\beginnumbering
\pstart
\setline{21}
Y el marcial Menelao, cuando lo vé que al frente\\
de las tropas, avanza con andar imponente,\\
se alegra como hambriento león que hace gran presa\\
de un enastado ciervo o una cabra montesa,\\
y ávido los devora, ni aunque esté perseguido\\
por los ágiles perros y los mozos fornidos.\par
\pend
%
\noindent-- -- -- -- -- --\par
%
\pstart
\setline{125}
Hallóla en el palacio tejiendo una amplia y bella\\
tela doble de púrpura en que historiaba entonces,\\
las mil luchas que a influjo de Ares sufren por ella\\
domadores troyanos y aqueos que arma el bronce.\par
\pend
%
\noindent-- -- -- -- -- --\par
%
\pstart
\setline{136}
El marcial Menelao y Alejandro, al instante,\\
por ti van a batirse con sus enormes lanzas,\\
y tú serás la esposa del que salga triunfante.\\
Diciendo así la diosa, con dulces esperanzas\\
de su primer marido, patria y deudos, la arroba.\\
Y al punto, envuelta en cándidos velos, deja la alcoba,\\
vertiendo tierno llanto...\par
\pend
%
\noindent-- -- -- -- -- --\par
%
\pstart
\setline{154}
Al ver Helena hacia ellos marchaba, con voz queda\\
y en aladas palabras dijeron: --Justo es para\\
los troyanos y aqueos padecer tantos males\\
por tal mujer; pues a las diosas inmortales,\\
ella, terriblemente, se parece de cara.\par
\pend
%
\noindent-- -- -- -- -- --\par
%
\pstart
\setline{164}
Tú no eres la culpable, sino los dioses, de esta\\
guerra de los aqueos, de plorable y funesta.\par
\pend
%
\noindent-- -- -- -- -- --\par
%
\pstart
\setline{173}
Dulce me habría sido la peor muerte, cuando\\
por seguir aquí a tu hijo, dejé el tálamo amable,\\
hermanos, dulce hija y amigas adorables.\\
Y de que así no fuera, consúmome llorando.\par
\pend
%
\noindent-- -- -- -- -- --\par
%
\pstart
\setline{178}
Es el poderosísimo Agamenón Atrida,\\
a la vez tan cumplido rey como buen guerrero:\\
cuñado de esta perra, por mi honra inmerecida.\par
\pend
%
\noindent-- -- -- -- -- --\par
%
\pstart
\setline{241}
\firstlinenum{242}
¿O no querrán, acaso, concurrir al combate,\\
corridos del oprobio que sobre mí se abate?\par
\pend
%
\noindent-- -- -- -- -- --\par
%
\pstart
\setline{320}
\firstlinenum{320}
Zeus padre que en el Ida reinas glorioso y grande:\\
haz que quien puso entre ambos pueblos el desconcierto,\\
descienda al fondo de la morada de Hades, muerto;\\
y que amistad, en cambio, tu alianza fiel nos mande.\par
\pend
%
\noindent-- -- -- -- -- --\par
%
\pstart
\setline{428}
Huiste de la guerra donde habría debido\\
inmolarte aquel héroe que antes fué mi marido.\\
A Menelao, amado de Ares, te lisonjeabas\\
de exceder en corage, brazo y lanza. Anda ahora\\
y rétalo de nuevo, como cuando peleabas...\\
Mas, nó, que disuadirte querría, previsora,\\
para que locamente no afrontes la pujanza\\
del blondo Menelao; no sea que éste al punto,\\
en la contienda bélica te rinda por la lanza.\par
\pend
%
\noindent-- -- -- -- -- --\par
%
\pstart
\setline{438}``Mujer, no me zahieras así. Si venció ahora\\
Menelao, con ayuda de Atena, vendrá mi hora;\\
porque también hay dioses que a nosotros nos guían.\\
Ea, pues, acostémonos y a querernos volvamos,\\
pues nunca embargó tanto tu amor el alma mía,\\
ni cuando de la amena Lacedemonia un día\\
te rapté, y en las naves, pasando el mar, bogamos,\\
y afecto y lecho uniéronnos en la lista de Kranae.\\
Así te amo y el ducle deseo a ti me atrae.\\
Tal dijo, y fuese al lecho, seguido de la esposa,\\
y ambos allá acostáronse en la cama lujosa''.
\pend
%
\noindent-- -- -- -- -- --\par
%
\endnumbering
\repage
%-
%Canto 4.
\head{\textsc{\emph{Ilíada}, Canto IV}}
\section*{\centering\textbf{\textsc{Canto} IV}\footnote{\ehdos, pp. 106-107.}}
\addcontentsline{toc}{subsection}{\textsc{Canto IV}}
%
\noindent-- -- -- -- -- --\par
%
\beginnumbering
\pstart
\firstlinenum{105}
\setline{104}Dice Atena y persuádelo; con que, al punto, el menguado,\\
toma el pulido arco sacado de un lascivo\\
cabrón montés, que un día cazara él emboscado,\\
hiriéndolo en el pecho con su tiro furtivo\\
al saltar de una roca donde quedó tumbado.\\
Sus cuernos que medían dieciseis palmos, fueron\\
labrados y ajustados por el maestro armero\\
que con remates de oro dejó el arco acabado.\\
Fíjalo él bien en tierra, lo tiende, y al instante\\
sus compañeros pónenle los escudos delante;\\
no acometa la brava gente aquea y le impida\\
herir a Menelao, marcial vástago Atrida.\\
Destapando él la aljaba, saca una nueva y pronta\\
saeta, mensajera del negro dolor; monta\\
en el nervio la amarga flecha, y una cumplida\\
hecatombe de blancos, primiciales corderos\\
prometiéndole a Apolo Licio, el ilustre arquero,\\
cuando vuelva en la sacra Zelea a ver su hogar--\\
tira juntos la muesca y el nervio, hasta acercar\\
la cuerda a su tetilla y el fierro a la otra parte\\
del corvo ingenio; y cuando lo acaba así de armar,\\
silba el arco, la cuerda brama, y aguda parte\\
la flecha hacia las filas, ávida de volar.\par
\pend
\endnumbering
%
\noindent-- -- -- -- -- --\par
\repage
%
%Canto 5.
\head{\textsc{\emph{Ilíada}, Canto V}}
\section*{\centering\textbf{\textsc{Canto} V}\footnote{\ehdos, pp. 113-131.}}
\addcontentsline{toc}{subsection}{\textsc{Canto V}}
%
\beginnumbering
\pstart
\firstlinenum{5}
\lettrine[lines=4,nindent=2pt]{P}{\textsc{alas Atena, entonces}}, dió audacia y entereza\\
a Diomedes Tideides, para que alta proeza\\
destacara entre todos los argivos su fama.\\
Le encendió escudo y yelmo con incansable llama,\\
tal como la del astro de Otoño, que lavado\\
por el Océano emerge más que todos brillante:\\
así el fuego en sus hombros y cabeza ha brotado,\\
cuando lo echa en lo recio del tumulto, adelante.\par
Había en Troya un cierto Dares, rico y honrado,\\
sacerdote de Hefesto, cuyos dos hijos, que eran\\
Ideo y Fegeo, aptos para todo combate,\\
separados del resto, contra el héroe se abaten\\
montando un carro, mientras él a pie los espera.\\
Fegeo, al encontrarse , tiró primeramente\\
su venablo de larga sombra que inútilmente\\
por sobre el hombro izquierdo del Tideida pasó.\\
Pero éste con segura mano le arrojó el bronce\\
que en la mitad del pecho fué a clavársele entonces,\\
tumbándolo del carro, pues no en vano partió.\\
Corrió Ideo, el hermoso vehículo dejando,\\
sin defender el cuerpo del hermano caído.\\
Que él mismo de la negra muerte no hubiera huído,\\
si en tinieblas Hefesto no lo envuelve, evitando\\
con salvarlo, que el viejo quede más afligido,\\
y el hijo de Tideo tomó y dió los corceles\\
a los suyos, haciéndolos llevar a los bajeles.\\
Al notar los troyanos que huye un hijo de Dares\\
y que el otro ante el carro muerto quedó, se apena\\
su corazón magnánimo; y la ojizarca Atena\\
así dijo, domándole la mano, al furioso Ares:\par
--Ares, Ares, funesto, destructor, sitiador:\\
¿no es mejor que dejemos que aqueos y troyanos\\
luchen, y que la gloria les dé Zeus soberano,\\
mientras así eludimos su divino rencor?\par
Tal dijo, y del combate sacando a Ares furioso,\\
lo hizo sentar a orillas del Escamandro herboso.\\
Los dánaos pusieron en fuga a los troyanos,\\
y cada jefe, entonces, mató un hombre a sus manos.\\
Agamenón, rey de hombres, va primero y derriba\\
del carro al grande Odíon, rey de los halizones,\\
que dió la espalda, hundiéndole su lanza desde arriba,\\
y el pecho atravesándole por entre los pulmones.\\
Tal cayó, y con estrépito sus armas resonaron.\par
Mató allá Idomeneo, por la lanza preclaro,\\
a Festo, hijo de Boro el meonio, que había\\
llegado de la fértil Tarne, hundiendo la enorme\\
asta en su hombro derecho, cuando al carro subía.\\
Cayó así; horribles sombras entonces lo rodearon,\\
y los hombres de Idomeneo lo despojaron.\par
Menelao el Atrida mató a Escamandrio, experto\\
cazador, que de Estrofio fué hijo, con su asta de haya.\\
La propia Artemis era quien le enseñó el acierto\\
en cazar cuantas fieras en bosques y cerros haya.\\
Mas de nada valiéronle ni artemis la flechera,\\
ni su celebrado arte de tirar al acecho,\\
pues el ilustre Atrida que huyendo ante él lo viera,\\
lo lanceó entre los hombros, traspasándole el pecho.\\
Y así, al caer de bruces, sus armas resonaron.\\
--Merión mató a Fereclón, hijo de Harmón maestro\\
que fuera en toda clase de artes manuales diestro,\\
pues Palas Atenea le daba su alto amparo.\\
Fué él quien hizo a Alejandro los navíos veloces\\
que tanto mal a Troya y a él mismo ocasionaron,\\
por no haber entendido su oráculo a los dioses.\\
Merión al alcanzarlo, pues sin cesar lo hostiga,\\
por la nalga derecha le envasó el bronce fuerte;\\
y como el hueso hendiérale, tocando la vejiga,\\
gimió al caer de hinojos y lo envolvió la muerte.\par
Meges mató a Pedeo, bastardo de Antenor,\\
a quien como hijo propio la divina Teano,\\
por ser grata a su esposo, criara con amor.\\
El ilustre Fileides que a él se hallaba cercano,\\
fué y le clavó en la nuca su aguda lanza entonces,\\
y la lengua cortándole, le atravesó los dientes:\\
que así rodó por tierra, mordiendo el frío bronce.\par
Eurípilo Evemónides da al divino Hipsenor\\
hijo de Dolopión el Grande, que oficiaba\\
el culdo de Escamandro, y a quien el pueblo honraba\\
como a un dios, fiero tajo cuando huye con pavor.\\
La espada, junto al hombro, tronchó el brazo pesado,\\
y rodó por los suelos el miembro ensangrentado.\\
Y cerraron los ojos de aquél en la cruenta\\
lid, la muerte purúrea y la Parca violenta.\par
Y mientras así agítalos el combate empeñoso,\\
no se sabe si el hijo de Tideo, en su brío,\\
combate con troyanos o aqueos, tan furioso\\
lánzase a la llanura cual rebalsado río\\
que durante el deshielo va arrollando bravío\\
los diques y los setos de los campos feraces,\\
y al crecer con las lluvias de Zeus, en su desvío,\\
muchas bellas labranzas de los mozos deshace:\\
tal las densas falanges troyanas descompone\\
Tideides, y entre tantos, ninguno se le opone.\\
Cuando, en eso, el caro hijo de Licaón, que advierte\\
que en la tropa del llano siembra, feroz, la muerte,\\
tiende el corvo arco y clávale en la espalda derecha,\\
por el hueco de la honda coraza, amarga flecha\\
que ensagrentando el bronce, del otro lado asoma.\\
Y el claro Licaónida grita con voz terrible:\par
--¡Cargad, bravos troyanos, hábiles en la doma!\\
Herido está el aqueo más bravo, y no es posible\\
que el grave tiro aguante, si con suerte propicia\\
me endilgó acá el rey, vástago de Zeus, desde la Licia.\par
Tal se alaba, aunque al otro no rindió el arma aguda;\\
pues ante yunta y carro reculó con pericia,\\
pidiendo a Esténelo, hijo de Capaneo ayuda:\par
--¡Corre, buen Capaneides, deja el carro, y prolijo\\
\dlgo{arráncame del hombro la amarga flecha!}{Dijo,}\\
y al punto apeóse Esténelo del carro, y la enemiga\\
flechada clavada al hombro le arrancó; y desbordando,\\
saltó la sangre por las mallas de la loriga.\\
Y así fue el buen guerrero Diomedes implorando:\par
--Escúchame hija indómita de Zeus potente: si\\
tu amparo en la ardua lucha nos diste siempre a mí\\
y a mi padre, ahora haciéndome la gracia, oh Atenea,\\
de que mate yo a ese hombre, ponlo a tiro de lanza;\\
pues porque me hirió, alábase con la loca esperanza\\
de que yo, pronto, el caro fulgor del sol no vea.\par
Así rogó y fué oído por Palas Atenea,\\
qe aligerando al punto su cuerpo, pies y manos,\\
con aladas palabras lo animó a la pelea:\par
--Recóbrate, Diomedes, y ataca a los troyanos,\\
que ya infundí en tu pecho la audacia y el vigor\\
de tu intrépido padre Tideo el domador,\\
hábil en el escudo; y aun torné caediza\\
la nube de ceguera que tapaba tus ojos,\\
para que a dioses y hombres conozcas en la liza,\\
y si un numen te prueba, no afrontes sus enojos.\\
Sólo entre Afrodita, la hija de Zeus, entonces\\
hiérela sin escrúpulo con el agudo bronce.\par
Dicho esto, fuése Atene la ojizarca; y al frente\\
de combate, Tideides se incorporó sin falta,\\
pues contra los troyanos, más aún que anteriormente,\\
triplicósele el brío, como al león que asalta\\
un establo de ovejas, donde lo ha lastimado\\
el pastor sin rendirlo, con que su furia exalta.\\
Húyese el hombre, entonces, al fondo del cercado;\\
y mientras indefenso y en su terror febril\\
con revuelto atropello se amontona el ganado,\\
la ávida fiera salta del profundo redil.\\
Así el fuerte Diomedes, en furia arrebatado,\\
las falanges troyanas revuelve y desbarata.\par
Allá a Astino y a Hipéiron, pastor de pueblos, mata,\\
clavando a éste la pica de bronce en la tetilla,\\
mientras con la ancha espada troncha al otro la istilla,\\
y de la espalda y cuello todo el hombro separa.\\
Déjalos allá y cierra contra Abante y Polido\\
a los que un viejo intérprete de sueños engendrara:\\
Eurídamas, que al verlos partir no habrá entendido\\
los suyos, pues Diomedes a ambos mata y despoja.\\
Acto continuo contra Janto y Toón se arroja,\\
hijos que pudo Fénope en su vejez lograr,\\
pues ya caduco y triste, no engendró otro heredero.\\
Al quitarles la amable vida, en llanto y pesar\\
sumió al padre, que vivos no ha de verlos tornar:\\
con los que otros parientes la herencia se repartieron.\par
Luego en un carro toma dos héroes que hijos fueron\\
de Príamo Dardánida: Cromíon y Equemón.\\
Cual desnuca, asaltando la boyada, el león,\\
a un buey o una ternera que en el bosque pacía,\\
Tideidos lo desmonta por la fuerza, con impía\\
saña, pilla sus armas, y entrega los corceles\\
a los suyos, haciéndolos llevar a los bajeles.\par
Mas, cuando destrozaba las filas, viólo Eneas,\\
y por entre el tumulto de armas de la pelea\\
fué en busca del divino Pándaro. Así que halló\\
al eximio y bravo hijo de Licáon, lo encaró\\
\dlgo{para decirle:}{--Pándaro, ¿qué es de tu arco y aladas}\\
flecha, y de tu fama por nadie disputada\\
aquí ni en Licia, donde ninguno se gloría\\
de aventajarte? ¡Vamos! levanta a Zeus tus manos,\\
y tira contra ese hombre, que haciendo a los troyanos\\
innumerables males, triunfa con saña impía,\\
y así a tanto valiente las rodillas afloja;\\
si no es algún dios que hacia los de Troya indispuso\\
un sacrilegio, porque malo es si un dios se enoja.\par
Entonces el claro hijo de Licáon le repuso:\\
--Eneas que a los nobles troyanos aconsejas,\\
en todo al bélico hijo de Tideo semeja.\\
Su escudo, su oculario triple casco ver creo,\\
y sus corceles, aunque puede que sea un dios;\\
mas, si es el bélico hijo del ilustre Tideo,\\
su furia está probando que un inmortal va en pos,\\
envuelto en una nube, y apartando la aguda\\
flecha que lo amenaza; pues ya una le tiré,\\
y por el hueco de la coraza le acerté\\
a la espalda derecha, creyendo que sin duda\\
al Hades lo echaría, pero no lo maté:\\
con que, seguramente, será un dios irritado.\\
Ni caballos ni carro tengo aquí qué montar,\\
aun cuando en el palacio de Licáon han quedado\\
once hermosos y sólidos coches sin estrenar,\\
con los toldos dispuestos, y cada uno a su lado\\
la yunta que cebada blanca y avena come.\\
El anciano guerrero Licáon prescribióme\\
al dejar yo el palacio, con premiosa insistencia,\\
que a los troyanos sólo desde el carro mandara\\
en la ardua lucha, pero no le presté aquiesencia.\\
Cuánto más me valiera, pues dejé por prudencia\\
los caballos, temiendo que el pasto les faltara\\
con reunión tan grande, y estando acostumbrados\\
a comer bien; sin ellos, a pie a Ilión he venido,\\
confiado sólo en mi arco, que nada útil me ha sido;\\
pues ya contra dos jefes, Tideides y un Atrida\\
he tirado, y aunque a ambos sangré con grave herida,\\
sólo excitarlos pude. Con mal sino partí\\
el día en que el corvo arco del clavo desprendí\\
y a la amena Ilión traje mis troyanos, por ser\\
amable al divino Héctor. Mas, como vuelva a ver\\
a mi patria, a mi esposa y a mi alta fortaleza,\\
que un guerrero enemigo me corte la cabeza,\\
si este arco no destrozo y echo al fuego violento,\\
ya que su compañía me es vana comp el viento.\par
Mas, replicóle el jefe de troyanos, Eneas:\\
--No hables así, pues nada cambiará, si nosotros\\
no vamos contra ese hombre con armas, carro y potros\\
a pelearlo. ¡Ea! al mío sube, para que veas\\
cómo saben lo mismo los corceles troyanos\\
perseguir o escaparse rápido por los llanos,\\
y cómo han de llevarnos salvos a la ciudad,\\
si otra vez Zeus concede gloria, en su potestad,\\
a Diomedes Tideides. Toma, pues, ya, a mi lado,\\
fusta y brillantes riendas; yo pelearé montado,\\
o hazlo tú, y los caballos a mi cuidado sean.\par
\dlgo{Y contestó el claro hijo de Licaón:}{--Eneas,}\\
encárgate tú mismo de la yunta y las bridas;\\
pues hallándose bajo su habitual conductor,\\
si el hijo de Tideo nos forzara a la huída,\\
tirará aquélla el cóncavo carro mucho mejor;\\
no se asuste y desboque, tu voz desconociendo,\\
y a sacarnos se niegue de la lucha, y cayendo\\
sobre ambos el bravo hijo de Tideo, nos mate\\
y se lleve los fuertes potros. Así, al combate\\
conduce por tu propia mano tu carro y yunta,\\
mientras su ataque espero con mi lanza de punta.\par
Hablando así, al labrado carro suben, y ansiosos,\\
contra Tideides lanzan los corceles briosos.\\
Mas Esténelo, el célebre hijo de Capaneo,\\
que los advierte, dícele con palabras aladas:\par
--Carísimo Diomedes Tideides, venir veo\\
dos robustos varones, con ardiente deseo\\
de pelearte, y que tienen fuerzas desmesuradas:\\
Pándaro el diestro arquero, quien se jacta de que es\\
Vástago de Licáon, y Eneas, que a su vez\\
se alaba de que Anquises, varón irreprochable,\\
lo engendró en Afrodita. Retirémonos, pues,\\
montados; no te expongas, furioso, a algún revés,\\
si en la vanguardia arriesgas perder la vida amable.\par
Mas el fuerte Diomedes, mirándolo de abajo,\\
\dlgo{contestó:}{--No me digas que huya o te des trabajo}\\
de convencerme en vano, pues no soy del linaje\\
de los que en retirada, temerosos pelean.\\
Iré a ellos desmontado, que aun me sobra coraje,\\
firme con el apoyo de Palas Atenea.\\
No han de llevarlos lejos de acá, si retroceden,\\
sus veloces caballos, y alguno escapar puede.\\
Mas, te diré algo, y guárdalo fielmente en tu memoria:\\
si la sapiente Atena me concede la gloria\\
de matarlos, tú nuestros rápidos troncos ata\\
al barandal, ciñéndoles las riendas, y arrebata\\
lejos de los troyanos los corceles de Eneas,\\
que pondrás al amparo de las tropas aqueas;\\
pues su raza es la de esos que a Tros dió en expiación\\
por Ganimedes, su hijo, Zeus el clarividente.\\
Nada hay, bajo la aurora y el sol, más excelente.\\
El rey de hombres Anquises, burló a Laomedón\\
y sacó cría, echándoles yeguas ocultamente.\\
Naciéronle seis de ese linaje en su mansión:\\
cuatro crió a pesebre, y a Eneas hizo don\\
de los otros dos, que ahora sembrar el terror vemos.\\
Con que, si los tomamos, un buen triunfo obtendremos.\par
Así éstos conversaban, mientras que los dos otros\\
caían ya sobre ellos con sus veloces potros.\\
Y este reto el claro hijo de Licáon les echa:\par
--¡Corajudo y noble hijo del ilustre Tideo,\\
pues no llegó a postrarte la amarga y rauda flecha,\\
veremos si la lanza responde a mi deseo!\par
Dice, y cimbrando tira su asta de larga sombra,\\
y el broquel del Tideida con el bote atraviesa,\\
rozando la coraza, y así a jactarse empieza:\par
--¡Atravesado tienes el vientre, y no me asombra\\
que pronto acabes, dándome considerable fama!\par
Pero el fuerte Diomedes, sin conmoverse, exclama:\\
--¡Erraste y falló el tiro! Mas creo por mi parte\\
que no lograréis tregua, como uno de ambos no harte\\
de sangre al bélico Ares, rodando por la arena.\par
Dice y tírale; el dardo, guiado por Atena,\\
entre ojo y nariz clávasele, sus blancos dientes pasa,\\
y el incansable bronce la lengua le cercena\\
de raíz, y su punta bajo el mentón rebasa.\\
Así cayó del carro; resonaron sobre él\\
sus magníficas armas, se abrieron sus corceles,\\
y allá lo abandonaron su fuerza y su alma fiel.\par
Eneas salta al suelo con su lanzón y escudo,\\
y temiendo que el cuerpo le quiten los aqueos,\\
ronda en torno, confiado como un león forzudo,\\
prontos broquel y pica, con ardientes deseos\\
de matar al que ataque, mientras grita sañudo.\\
Tideides alza entonces una piedra imponente\\
que dos hombres de ahora levantar no podrían,\\
mientras con una mano la lleva él fácilmente.\\
Con ella hiere a Eneas donde el fémur encaja\\
con la cadera, punto que se llama \emph{cotylo};\\
rómpele ambos tendones y el \emph{cotylo} descuaja,\\
y el cutis le desgarra la piedra con su filo.\\
El héroe, de rodillas cae, apoyando en tierra\\
su fuerte mano, y negra noche sus ojos cierra.\par
Y allá Eneas, rey de hombres, habría perecido,\\
si acaso Afrodita, hija de Zeus, que lo tuviera\\
de Anquises al boyero, pronto no lo advirtiera.\\
Con que, sus blancos brazos tiende al hijo querido,\\
y en un doblez del mando brillante lo resguarda,\\
no sea que los dánaos bien montados le claven\\
el cruel bronce en el pecho, y así su vida acaben.\par
Y mientras del combate lo substrae, no tarda\\
en efectuar Esténelo lo que convino, al mando\\
del cumplido guerrero Diomedes. Apartando\\
la yunta del tumulto, sus bridas ciñe al tope\\
del barandal, y leos de los troyanos lleva\\
los lucientes caballos de Eneas, al galope,\\
rumbo hacia los aqueos de las hermosas grebas\\
los da a Deipilo (amigo que entre todos honraba,\\
por lo bien que su juicio con el suyo acordaba)\\
a fin de que a los buques los conduzca; y montando\\
su yunta otra vez, coge más bien labradas riendas,\\
y solícito anímala, a Tideides buscando.\\
Este acosa a Ciprina con sus armas tremendas,\\
sabiéndola una débil diosa entre las deidades,\\
no de esas que a los hombres guían en las contiendas,\\
como Atena o Enío que arrasa las ciudades.\\
Y así que en la refriega tumultuosa la alcanza,\\
el hijo de Tideo cala su aguda lanza\\
y con ella atacándola, rompe el mando divino,\\
propia obra de las Gracias, y herir el cutis fino\\
de la mano, al extremo del dorso palmar, osa.\\
Brota entonces la sangre divina de la diosa,\\
o mejor dicho el ícor, flúido peregrino\\
que los dioses dichosos en vez de ella derraman,\\
porque ni trigo comen, ni beben negro vino:\\
(con lo que, siendo exangües, inmortales los llaman).\\
Lanza ella agudo grito, dejando a su hijo solo;\\
y mientras lo alza en una nube azul Febo Apolo,\\
no sea que los dánaos bien montados le claven\\
el cruel bronce en el pecho, y así su vida acaben,\\
increpa el buen guerrero Diomedes a la diosa:\par
--¡Hija de Zeus, apártate de la acción belicosa!\\
¡Basta ya con que engañes a débiles mujeres!\\
Pues creo que si en otra guerra mezclarte quieres,\\
hasta su mismo nombre te alejará medrosa.\par
Dijo, y ella apartóse cruelmente torturada.\\
La rauda Iris sacóla del tumulto, abrumada,\\
negreándole ya el lindo cutis. Pronto halló al fiero\\
Ares, que hacia la izquierda del combate sentado,\\
contra una nube había su lanzón arrimado,\\
lo propio que la yunta de caballos ligeros.\\
Y ante el querido hermano se arrodilló, abatida,\\
instándole a prestarle los potros de áurea brida:\par
--Querido hermano, apiádate de mi y préstame ahora,\\
para ir al Olimpo en que los inmortales moran,\\
tus caballos; pues sufro mucho de hallarme herida\\
por un varón, Tideides, quien con Dios padre fuera\\
\dlgo{capaz de Pelear.}{Díjole así, y Ares le ha dado}\\
los potros de áurea brida, con que el carro ha montado,\\
gimiendo de congoja. Subió Iris a su vera,\\
y picó, rienda en mano, la yunta voladora.\\
Pronto al Olimpo fueron, donde los dioses moran;\\
y la diligente Iris, los potros sujetando,\\
los desató y echóles por pasto la ambrosía.\par
La divina Afrodita, mientras tanto, caía\\
al regazo de Dione su madre, que tomando\\
en los brazos a su hija, la acarició y le dijo:\par
--¿Qué uraniano, hija amada, te hizo esta iniquidad,\\
como si procedieras con notoria maldad?\par
\dlgo{La risueña Afrodita respondióle:}{--Fué el hijo}\\
de Tideo, el soberbio Diomedes, quien me hirió,\\
porque a Eneas, que es mi hijo más caro, saqué yo\\
del combate. Pues ahora ya no son entre iguales,\\
aqueos y troyanos, las luchas más furiosas;\\
que los dánaos atacan hasta a los inmortales.\par
Y le contestó Dione, divina entre las diosas:\par
--Soporta, hija, tus penas, que muchos moradores\\
del Olimpo, debimos sufrir tormento agudo\\
de los hombres, mezclándonos así con sus rencores.\\
Los pasó Ares, cuando Oto y Efialtes el forzudo,\\
hijos de Alo, apresáronlo con formidable nudo,\\
por trece meses dentro de un cántaro de bronce.\\
Y muriera allá el bélico Ares, si Eríbea, entonces,\\
bellísima madrastra de aquéllos, no lo advierte\\
a Hermes, quien substrajo a Ares del lazo rudo y fuerte\\
que lo postraba exánime. Hera los padeció,\\
cuando el vigoroso hijo de Anfitrión le clavó\\
en el seno derecho saeta trifurcada,\\
en tormento incurable dejándola angustiada.\\
Y sufrió el enorme Hades la fecha audaz de aquel\\
mismo hombre, hijo de Zeus Portaégida, cuando ante\\
la puerta de los muertos le infligió herida cruel.\\
De dolor traspasado, y el corazón penante,\\
pues abatía su ánimo la saeta enterrada\\
en su espalda robusta, fué al Olimpo, morada\\
de Zeus, donde aplicándole Peón drogas calmantes,\\
lo curó, pues no había nada en él de mortal.\\
¡Bárbaro y temerario quien con impiedad tal,\\
por el arco a los dioses del Olimpo abatió!\\
Así Atena, la zarca diosa, contra tí alzó\\
al hijo de Tideides, quien, insensato, ignora\\
que todo el que a los dioses combate en mala hora,\\
no alcanza larga vida, ni logra que ocupando\\
sus rodillas los hijos, le llamen padre, cuando\\
vuelve de los combates y la feroz pelea.\\
Piense Tideides, aunque se tenga por tan fuerte,\\
que alguien puede atacarlo mejor que tú; no sea\\
que muy luego la próvida Adrastina Egialea,\\
corte, gimiendo, el sueño, y en su inquietud despierte\\
a las caras doncellas, llorando a su marido\\
legítimo, ese aqueo principal y cumplido,\\
el domador Diomedes, de quien es noble esposa.\par
Dice y resta el ícor, en las suyas tomando\\
la mano que así cura, su cruel dolor calmando.\\
Mas, como Atena y Hera presenciaran la cosa,\\
picaron con palabra mordaz al dios Kronida.\\
Y dijo, adelantándose, así la zarca diosa:\par
--Zeus padre, aunque te irrite mi palabra atrevida,\\
seguro es que Ciprina, tentando a alguna aquea\\
de hermoso peplo, para que siga a los troyanos\\
que tanto ama, al hacerle las caricias que emplea,\\
se arañó en áureo bronce la delicada mano.\par
Tal dijo. El padre de hombres y dioses sonrió,\\
y a la áurea Afrodita llamando, así le habló:\par
--El bélico trabajo no es para tí, hija mía,\\
sino los de himeneo, con su dulce alegría.\\
Que cuiden de eso Atena y Ares el dios ligero.\par
Mientras así conversan, lánzase el buen guerrero\\
Diomedes contra Eneas, aunque sabe que Apolo\\
le tiene de su mano, pues ni al gran dios flechero\\
respeta ya, deseando matar a aquél tan sólo,\\
y quitarle el magnífico arnés. Tres veces carga\\
con ansia de ultimarlo, y otras tantas Apolo\\
se le opone, agitando su refulgente adarga.\\
Mas, cuando igual a un numen por cuarta vez se obstina,\\
Apolo, el grande arquero, terrible lo amenaza:\par
--Tidida, piensa y vete; no quieras darte traza\\
de igualarte a los dioses, que la raza divina\\
nunca se asemejó a la terrestre humana raza.\par
Así dice, y Tideides recula un poco, sólo\\
por evitar la cólera del flechador Apolo.\par
\pend
%
\noindent-- -- -- -- -- --\par
%
\pstart
\setline{584}Antíloco, atacándolo, le hundió en la sien la espada,\\
con que, del carro hermoso, cayó al suelo, anhelante,\\
quedando un largo rato de cabeza, clavada\\
toda ella hasta los hombros en la arena abundante.\\
Tendiólo al fin por tierra la yunta alborotada,\\
que picó hacia el ejército aqueo el valeroso\\
Antíloco, Héctor, viéndolos, se arroja clamoroso,\\
con las recias falanges troyanas que comanda\\
Ares, al propio tiempo que Enío veneranda.\\
Va ésta con el Tumulto feroz y escandaloso,\\
mientras el dios, blandiendo su lanza prepotente,\\
ya marcha a espaldas de Héctor, ya anda con él al frente.\\
Tembló al verlo Diomedes el bravo, a la manera\\
del que yendo extraviado por una gran pradera,\\
dá con rápido río que echa al mar su corriente,\\
y al ver la hirviente espuma, salta retrocediendo,\\
así paró Tideides, a los suyos diciendo:\par
--Oh amigos ¿cómo habremos de asombrarnos que sea\\
Tan bravo el divino Héctor y audaz en la pelea?\\
Siempre con él va un numen que la muerte le evita,\\
Y ahora es Ares, que a un hombre mortal junto a él imita.
\pend
\separator
%
\pstart
\setline{720}
Aprontar los corceles de áurea brida, dispuso\\
la augusta diosa Hera, hija del gran Kronos, entonces.\\
A su vez la pronta Hebe fué y al coche le puso\\
en el eje de hierro combas ruedas de bronce,\\
de ocho rayos, que sobre pinas de oro inusable\\
llantas de bronce ostentan con ajuste admirable;\\
mientras que son de plata los cubos bien torneados.\\
Y de oro y plata tiene la caja sus sopandas,\\
y por la delantera le corren dos barandas.\\
El timón es de plata, y ella a su extremo ha atado\\
yugo y coyundas de oro de gran belleza. Y hera,\\
avida de combates y ejercicio agitado,\\
unce al yugo los potros firmes en la carrera.\par
Mientras tanto Atena, hija del dios portaéjida, echa\\
sobre el umbral paterno su velo bien bordado\\
que sus manos tejieron, y luego se ha ajustado\\
el arnés del nubígero dios: que así se pertrecha\\
para la guerra aciaga; y a la espalda ha cargado\\
la fiera y floqueada égida que el \emph{Espanto} corona.\\
En ella están la \emph{Fuerza}, la \emph{Discordia}, y la atroz\\
\emph{Persecución}, y la hórrida testa de la \emph{Gorgona},\\
prodigio de Zeus, monstruo tremebundo y feroz.\\
Armaste la cabeza con una áurea celada\\
de doble penacho y de cuádruple carrillera,\\
que a los infantes juntos de cien pueblos cubriera;\\
y montando el fulgente carro, ase la pesada\\
ingente y recia lanza con que a tanto valiente\\
rinde, si en contra suya se irrita prepotente.\par
\pend
\separator
\endnumbering
\repage
%
%Canto 6.
\head{\emph{\textsc{Ilíada}}, Canto VI}
\section*{\centering\textbf{\textsc{Canto} VI}\footnote{\ehuno, pp. 42-43 + \ehdos, pp. 86, 94-97 + \ehctr, pp. 262-265.}}
\addcontentsline{toc}{subsection}{\textsc{Canto VI}}
%
\beginnumbering
\noindent-- -- -- -- -- --\par
\pstart
\setline{123}
¿Quién de entre los mortales eres tú el esforzado?\\
Jamás te ví en la lucha que a los héroes renombra;\\
mas, por tu audacia a todos los has sobrepasado,\\
con afrontar mi lanza, la de la larga sombra.\par
\pend
\noindent-- -- -- -- -- --\par
%
\pstart
\setline{128}
\firstlinenum{129}
Mas, si del cielo vienes, por ser un inmortal,\\
no pelearé, por cierto, contra un Dios celestial.\par
\pend
\noindent-- -- -- -- -- --\par
%
\pstart
\setline{142}
\firstlinenum{143}
Mas, si un mortal que nutre la tierra, eres por suerte,\\
acércate y más pronto te ultimará la muerte.\par
\pend
\noindent-- -- -- -- -- --\par
%
\pstart
\setline{145}
\firstlinenum{145}
Magnánimo Tideides, ¿a qué indagas mi origen,\\
si el del hombre es como el de las hojas pasajeras\\
que el viento en tierra esparce, cuando ya las otras se erigen\\
al pintar en el fértil bosque de la primavera?\\
Tal las generaciones nacen y acaban; pero\\
ya que saber desear nuestro bien conocido\\
linaje, que te enteres cumplidamente quiero.\par
\pend
\noindent-- -- -- -- -- --\par
%
\pstart
\setline{212}
Tal se expresó. Diomedes el ilustre en la guerra,\\
contento de escucharlo, clavó su lanza en tierra,\\
y así al pastor de pueblos dijo con voz amable:\par
--Eres mi antiguo huésped paterno, en verdad, pues\\
el deífico Eneo dió posada una vez\\
por veinte días en su palacio al intachable\\
Belerofonte, y ambos se obsequiaron cumplidos:\\
dió Eneo un cinto espléndido, de púrpura teñido,\\
y Belerofonte una copa de oro con asas,\\
que cuando yo me vine para acá, quedó en casa.\\
No recuerdo\footnote{En el impreso figura como «recuerbo», probablemente un error de imprenta.} a Tideo, puesto que me dejó\\
muy niño, cuando en Tebas pereció el pueblo aqueo;\\
mas, tu amistoso huesped ser en Argos deseo,\\
cual lo serás tú en Licia si a ese país voy yo.\\
Evitemos, pues, nuestras lanzas en la refriega,\\
que a mí muchos troyanos y aliados qué matar\\
me quedan todavía, si un dios me los allega,\\
y a tí muchos aqueos que podrás ultimar.\\
Y ahora cambiemos armas, y así éstos podrán ver\\
que huéspedes paternos nos gloriamos de ser.\par
\pend
\separator
%
\pstart
\setline{232}
Dicho esto, y de los carros descendiendo en seguida,\\
danse las manos para que su amistad se pruebe.\\
Entonces turbó a Glauco la razón Zeus Kronida;\\
pues al cambiar sus armas con Diomedes Tidida,\\
dió aquél oro por bronce, cien bueyes contra nueve.\par
\pend
\separator
%
\pstart
\setline{274}
Y prométele doce terneras, que inmoladas\\
de un año, y aun no uncidas, le serán, si se apiada\\
de la ciudad, esposas y chiquillos troyanos,\\
y al hijo de Tideo, guerrero que inhumano\\
nos amedrenta, aleja nuestra Ilión sagrada.\par
\pend
\noindent-- -- -- -- -- --\par
%
\pstart
\setline{305}
Divina entre las diosas, venerable Atenea,\\
Patrona nuestra, rompe la lanza de Diomedes\\
y haz que caiga delante de las Puertas Esceas.\par
\pend
\noindent-- -- -- -- -- --\par
%
\pstart
\setline{344}
Oh hermano de esta perra funesta y horrorosa,\\
¿por qué cuando mi madre me dió a luz, un mal viento\\
no me arrojó a los montes o al mar que, violento,\\
me habría arrastrado, antes de ocurrir estas cosas?\\
Pero, ya que los dioses tenían el intento\\
de causar tantos males, yo debí ser la esposa\\
de un hombre mejor, que ante la condena y la afrenta\\
de los otros varones, tomara su honra en cuenta.\\
De firmeza y corage siempre éste ha carecido,\\
hasta que un día, creo, tendrá su merecido.\\
Mas, entra hermano, y siéntate aquí, y reposa en calma,\\
del afán que por esta perra te angustia el alma,\\
y a causa de Alejandro; pues fatal anatema\\
nos ha infligido Zeus, para que en adelante,\\
sirva a los venideros nuestro baldón de tema.\par
\pend
\separator
%
\pstart
\setline{390}
Dijo el ama. Hérctor, yéndose de la casa muy presto,\\
por las bien hechas calles desanduvo el camino.\\
Ya en las puertas Esceas (que es de donde, al lado opuesto,\\
debía sobre el llano salir) corriendo vino\\
hacia él su esposa Andrómaca, opulenta heredera\\
del magnánimo Eécion que a la falda viviera\\
del Placo umbroso, en Tebas la hipoplaciana, donde\\
fué rey de los cilicios. A su hija poseía,\\
pues, Héctor del broncíneo casco. Así ella venía\\
ante él, y a la par suya marchaba una doncella,\\
llevando al seno el hijo de Héctor, un tierno infante,\\
que era su bienamado, lindo cual una estrella.\\
Escamandrio llamábalo él; mas los otros, como\\
sólo Héctor defendía la ciudad con aplomo,\\
Astianax le decían. El héroe, con cariño,\\
sonriendo en silencio contemplaba a su niño.\\
Andrómaca, llorosa, detúvose a su lado,\\
\dlgo{y de su mano asiéndole, díjole:}{--Infortunado,}\\
tu valor va a perderte, y en tanto no te apiadas\\
ni de tu hijo chiquito, ni de mí, desdichada,\\
que seré tu viuda pronto, pues sobre ti\\
caerán los aqueos para inmolarte así.\\
Privada de tu apoyo, morir será mi anhelo,\\
que tu muerte dejárame un dolor sin consuelo.\\
Ya mi padre no existe, mi augusta madre es muerta.\\
mató a mi padre Aquiles, asoló la poblada\\
ciudad de los cilicios, Tebas la de altas puertas,\\
mas no despojó a Fécion; antes bien, compungido,\\
quemólo con sus armas y un túmulo ha erigido,\\
que de álamos cercaron las ninfas Orcades\\
hijas de Zeus portaégida. Mis siete hermanos juntos,\\
en un día bajaron a las sombras del Hades,\\
que Aquiles, el de raudos pies, los dejó difuntos\\
entre los chuecos bueyes y las blancas ovejas.\\
Mi madre que a la falda de Placo nemoroso\\
reinaba también, trájola aquél, sordo a sus quejas,\\
con el botín. Mas, luego que mediante un cuantioso\\
rescate, concediónos que libertada fuera,\\
la inmoló en el alcázar Artemis la flechera.\\
Hoy, Héctor, para mí eres padre y madre honorable,\\
y hermano, a más de esposo, como ninguno amable.\\
Ten piedad, y no dejes la torre sin tu ayuda.\\
No hagas a tu hijo huérfano y a tu mujer viuda.\\
Detén a los soldados cerca del cabrahigo,\\
que ahí la ciudad es débil y el bastión inseguro.\\
Ya los héroes más bravos que cuenta el enemigo,\\
tres veces intentaron allá expugnar el muro,\\
bajo los dos Ayaces, el claro Idomeneo,\\
los Atridas y el fuerte vástago de Tideo.\\
O alguien quen los oráculos sabe se lo ha indicado\\
o bien su propio arrojo los ha precipitado.\\
Respondióle el Grande Héctor del casco tremolante:\par
--Mujer, bien sé todo eso, pero tambén me aterra\\
lo que dirán, si tímido me aparto de la guerra,\\
troyanos y troyanas de peplo rozagante.\\
Ni está en mi mano hacerlo\footnote{En la impresión «hacrelo», otro error.}, que siempre en la porfía\\
fuí valeroso y supe mantenerme adelante,\\
sosteniendo la gloria de mi padre y la mía.\\
Bien lo tengo en mi mente y en mi pecho arraigado:\\
día vendrá en que, al último, perezca Ilión la santa,\\
y Príamo y su pueblo, con el buen fresno armado.\\
Mas, ni el dolor de Troya desvalida me espanta,\\
ni el de la misma Hécuba, ni el de Príamo augusto,\\
ni el de mis muchos bravos hermanos, que el desastre\\
dará por tierra como tu dolor, cuando entonces,\\
uno de esos aqueos tunicados de bronce,\\
sollozante y privada de libertad te arrastre.\\
O cuando en Argos tejas para otro, y compelida\\
a ir por agua a las fuentes Meseida o Hiperea,\\
en ti sea la dura fatalidad cumplida.\\
Y quizá exclame alguno que llorando te vea:\\
esta es la esposa de Héctor que fué el primer guerrero\\
entre aquellos troyanos que ante Ilión combatieron.\\
Así dirán, y nuevo dolor ha de angustiarte\\
de haber perdido al hombre capaz de liberarte.\\
Pero antes cubra el polvo mis mortales despojos,\\
que oiga tu grito o vea tu rapto con mis ojos.\\
Dice y tiende sus brazos al niño; mas, gritando,\\
estréchase éste al seno de la apuesta nodriza,\\
pues lo asusta su tierno padre al irse allegando,\\
con el bronce y la mata de crin que tremolando\\
en el timbre del yermo, feramente se riza.\\
Padre y madre sonríen, y el bravo Héctor se quta\\
el centelleante casco que en tierra deposita.\\
Besa luego al infante y en sus brazos lo mece,\\
y a Zeus y demás dioses dirige así sus preces:\par
--Oh, Zeus y demás dioses, haced que mi hijo un día,\\
cual yo entre los troyanos, sea fuerte y glorioso.\\
Que sobre Ilión se vea reinando poderoso,\\
y en los comates digan de él por su bizarría:\\
``Pero éste vale mucho más que su padre!" Cuando\\
vuelva con los sangrientos despojos de un guerrero,\\
el alma de su buena madre regocijando.\par
Dice así y en los brazos de la querida esposa\\
depone al pequeñuelo, que ella en su perfumado\\
seno acoge, sonriendo todavía llorosa.\\
Entonces él, notándolo, la acaricia apiadado,\\
\dlgo{y dícele:}{--Oh, mi triste, no temas por mí, o creas}\\
que alguien me arroje\footnote{En el impreso «arorje».} al Hades, contrariando el destino,\\
pues no hay mortal que, tímido o bravo, eluda el sino\\
con que nació. Ya a casa vuelve y a tus tareas\\
de telar y del huso, ya  ver que las criadas\\
hagan su parte, mientras del trabajo guerrero\\
se encargan los soldados de Ilión, y yo el primero.\par
\pend
\endnumbering
\repage
%
%
\end{document}